\item (a) Se define la integral impropia (en el plano completo $\mathbb{R}^2$)
\[ I = \iint_{\mathbb{R}^2} e^{-(x^2+y^2)} \, dA = \int_{-\infty}^\infty \int_{-\infty}^\infty e^{-(x^2+y^2)} \, dy \, dx = \lim_{a \to \infty} \iint_{D_a} e^{-(x^2+y^2)} \, dA \]
en donde $D_a$ es el disco de radio $a$ y centro en el origen. Demuestre que
\[ \int_{-\infty}^\infty \int_{-\infty}^\infty e^{-(x^2+y^2)} \, dA = \pi \]
(b) Una definición equivalente de la integral impropia de la parte (a) es
\[ \iint_{\mathbb{R}^2} e^{-(x^2+y^2)} \, dA = \lim_{a \to \infty} \iint_{S_a} e^{-(x^2+y^2)} \, dA \]
en donde $S_a$ es el cuadrado con vértices $(\pm a, \pm a)$. Use esto para demostrar que
\[ \int_{-\infty}^\infty e^{-x^2} \, dx \int_{-\infty}^\infty e^{-y^2} \, dy = \pi \]
(c) Deduzca que
\[ \int_{-\infty}^\infty e^{-x^2} \, dx = \sqrt{\pi} \]
(d) Efectuando el cambio de variable $t = \sqrt{2}x$, demuestre que
\[ \int_{-\infty}^\infty e^{-x^2/2} \, dx = \sqrt{2\pi} \]
(Este es un resultado fundamental de la probabilidad y estadística).
